\section{$\boldsymbol{D}\mathbf{+1}$ 维中的场论}\label{sec:fieldtheory}

上一节得到的、随时间变化的规范场关联函数的费曼规则，
正是 $D+1$ 维中一个定域场论的费曼规则 \cite{ZinnZwanziger}。
这一点可以通过多少有些形式化的泛函积分操作来证明，
但也可以采取纯代数的观点，此时 $D+1$ 维中的作用量
仅仅是作为费曼规则的生成函数而出现。

\subsection{作用量}

除了边界条件 \eqref{eq:2.3} 之外，现在把场 $B_{\mu}(t,x)$
视为独立场。还需要引入一个具有纯虚分量的拉格朗日乘子场
$L_{\mu}(t,x)=L_{\mu}^a(t,x)T^a$。
$D+1$ 维中理论的作用量于是为 \cite{ZinnZwanziger}
\begin{align}
   \Stot=S+\Sgf+\Sgh+\Sfl,
   \label{eq:4.1}\\[2.5pt]
   \Sfl=-2\int_0^{\infty}\rmd t\int\rmd^Dx\,
   \tr\bigl\{L_{\mu}(t,x)
   \bigl(\partial_tB_{\mu}-
         D_{\nu}G_{\nu\mu}-\alpha_0D_{\mu}\partial_{\nu}B_{\nu}
   \bigr)(t,x)\bigr\}.
   \label{eq:4.2}
\end{align}
在这一框架下，流方程 \eqref{eq:2.4} 恰好就是把作用量对拉格朗日乘子场
变分所得的场方程。注意，后者不需要满足任何特定的边界条件。

\subsection{传播子}\label{ssec:propagators4}

作用量的二次型部分是
\begin{equation}
   \int_0^{\infty}\rmd t\int\rmd^Dx\, L^a_{\mu}(t,x)
   \bigl(
   \partial_tB^a_{\mu}-\partial_{\nu}\partial_{\nu}B^a_{\mu}
   -(\alpha_0-1)\partial_{\mu}\partial_{\nu}B^a_{\nu}
   \bigr)(t,x)
   \label{eq:4.3}
\end{equation}
与 $D$ 维中作用量二次型部分之和。体场 $B_{\mu}$ 与 $L_{\mu}$
因此只彼此耦合，但通过边界条件 \eqref{eq:2.3}
还存在与基本规范场的隐式耦合。这个复杂因素很容易克服：
只需代入
\begin{equation}
   B_{\mu}(t,x)=\int\rmd^Dy\,K_t(x-y)_{\mu\nu}A_{\nu}(y)+b_{\mu}(t,x).
   \label{eq:4.4}
\end{equation}
于是场 $b_{\mu}$ 满足齐次边界条件，
而式 \eqref{eq:4.4} 右端第一项是线性化流方程的解，
从而在作用量 \eqref{eq:4.3} 中不出现。

由这些论述可知：基本规范场与鬼场的传播子
与上一节给出的表达式一致。所有其他两点函数均为零，除了
\begin{equation}
   \left.\bigl\langle
   b^a_{\mu}(t,x)L^b_{\nu}(s,y)
   \bigr\rangle\right|_{\rm leading\;order}
   =\delta^{ab}H(t,x;s,y)_{\mu\nu},
   \label{eq:4.5}
\end{equation}
它由场方程
\begin{equation}
   \bigl\{\delta_{\mu\rho}\partial_t-
   \delta_{\mu\rho}\partial_{\sigma}\partial_{\sigma}
   -(\alpha_0-1)\partial_{\mu}\partial_{\rho}\bigr\}
   H(t,x;s,y)_{\rho\nu}=\delta_{\mu\nu}\delta(t-s)\delta(x-y)
   \label{eq:4.6}
\end{equation}
及边界条件 $\left.H_{\mu\nu}(t,x;s,y)\right|_{t=0,s>0}=0$
确定。这组方程的唯一解是
\begin{equation}
   H(t,x;s,y)_{\mu\nu}=\theta(t-s)K_{t-s}(x-y)_{\mu\nu},
   \label{eq:4.7}
\end{equation}
因此 $bL$ 传播子恰与流传播子重合。

最后，涉及 $B$ 场的两点函数可以通过把式 \eqref{eq:4.4}
与迄今所得结果结合起来算出。由于 $AL$ 与 $bb$ 传播子为零，
很快便知 $BL$ 传播子等于流传播子，
而 $BB$ 与 $AB$ 传播子等于上一节给出的表达式。

\subsection{顶点}

由作用量 \eqref{eq:4.1} 描述的理论，
除 $D$ 维理论的顶点外还具有来自体作用量 $\Sfl$ 的顶点。
回顾式 \eqref{eq:2.6} 与 \eqref{eq:2.13}，容易证明后者相互作用部分
\begin{equation}
   \left.\Sfl\right|_{\rm interaction}=
   2\int_0^{\infty}\rmd t\int\rmd^Dx\,
   \tr\bigl\{L_{\mu}(t,x)R_{\mu}(t,x)\bigr\}
   \label{eq:4.8}
\end{equation}
生成的 $LB^2$ 与 $LB^3$ 顶点恰好就是顶点 $\vbulk{2,0}$ 与 $\vbulk{3,0}$。

这样，前面各节费曼规则的所有顶点与传播子都被重新得到。
这些规则的一个略显不寻常之处在于存在一个场（拉格朗日乘子），
它只通过与其它场的混合而传播。
不过，费曼规则的代数结构完全标准，
除了把作用量按场幂次展开所导出的那些规定之外不涉及任何其他规定。

\begin{figure}[t]
\centerline{\includegraphics[width=6.5cm]{images/diagram8}}
\caption{带流线圈的图的示例。鉴于流传播子的推迟性质
以及维数正规化的规则，所有这类图都为零 \cite{ZinnZwanziger}。}
\label{fig:loops}
\end{figure}

\subsection{流线与流线圈}\label{ssec:flowlines}

流线代表 $BL$ 传播子。它们可以从流顶点出发并终止于流顶点，
但决不会连接到普通顶点。朝外与朝内的外流线分别是外 $B$ 线与外 $L$ 线。

带闭合流线的图虽然符合这些规则，
但若要与前一节的费曼规则完全匹配，
这样的图就应当不出现（参见第 \ref{sec:perturbation} 节末尾的要点 (d)）。
事实上这类图等于零 \cite{ZinnZwanziger}。
例如图~\ref{fig:loops} 中的图 1 与图 2 为零，
因为维数正规化使动量积分为零。若圈中有两个或更多个顶点
（如图 3），则时间积分为零，因为流传播子是推迟的，
从而迫使各顶点处的流时间被压缩到一个测度为零的范围
（鉴于动量积分已被正规化，时间坐标中的奇点不会出现）。

如果采用格点正规化，图 2 的消失可能得不到保证，
但总可以在作用量中加入一个鬼场，在代数层面
（即在费曼被积函数的层面上）消去流线圈 \cite{Zinn}。
不过，采用维数正规化时无需这一手段，故此处略去。
